\documentclass{article}

% Required for inserting images
\usepackage{graphicx}

% Required for inserting links
\usepackage{hyperref}

% Required for math equations
\usepackage{amsmath}  
\usepackage{stix}

% General data
\title{Formulario Analisi 1}
\author{\href{https://github.com/pietropeerani}{Pietro Peerani}}
\date{2024}

% Page formatting
\usepackage[a4paper, margin=2.5cm]{geometry}
\usepackage{fancyhdr}
\pagestyle{fancy}
\fancyfoot[L]{Pietro Peerani}
\fancyfoot[R]{Analisi 1}



\begin{document}
    
    \maketitle
    \newpage
    \thispagestyle{plain} % remove fancy style from this page
    \tableofcontents
    \newpage
    
    \section{Forme utili}
    \subsection{Bernoulli (Jacob)}
    \begin{equation}
        (1+x)^n > 1+nx
    \end{equation}
    
    \subsection{Numero di Nepero}
    \begin{equation}
        \lim_{n\rightarrow\infty}\biggl(1+\frac{1}{n}\biggl)^n
    \end{equation}
    \begin{equation}
        \lim_{n\rightarrow\infty}\biggl(1-\frac{1}{n}\biggl)^{-n}
    \end{equation}
    

    \section{Limiti di successioni}
\subsection{Rapporto tra polinomi}
\begin{equation}
    a_n = \frac{P_1(x)}{P_2(x)} \Longrightarrow \exists\lim_{n\rightarrow\infty} a_n = 
    \begin{cases}
        0 & se\ g_2 > g_1\\
        \frac{c_1}{c_2} & se\ g_2 = g_1\\
        \pm\infty & se\ g_2 < g_1
    \end{cases}
\end{equation}

\subsection{Teorema (delle sottosuccessioni)}
Supponiamo di avere $\{a_n\}$ tale che $\exists\lim_{n\rightarrow\infty}a_n$ (finito o infinito), allora per ogni sottosuccessione $\{a_{n_k}\}$ di $\{a_n\}$ si ha $\lim_{n\rightarrow\infty}a_{n_k} = \lim_{n\rightarrow\infty}a_n$

\subsubsection{Sottosuccessioni}
Se esistono 2 sottosuccessioni $\{a_{n_{k_n^1}}\}$ e $\{a_{n_{k_n^2}}\}$ tale che $\lim_{n\rightarrow\infty}a_{n_{k_n^1}} \not = \lim_{n\rightarrow\infty}a_{n_{k_n^2}}$. Allora (per il teorema precedente):
\begin{equation}
    \not\exists\lim_{n\rightarrow\infty}a_n
\end{equation}

\subsection{Criterio del rapporto}
Sia $\{a_n\} > 0$ tale che:
\begin{equation}
    \exists\lim_{n\rightarrow\infty}\frac{a_{n+1}}{a_n} = L
\end{equation}
Allora:
\begin{equation}
    \exists\lim_{n\rightarrow\infty} a_n = 
    \begin{cases}
        L \in [0, 1) & \rightarrow \exists\lim_{n\rightarrow\infty} a_n = 0 \\ 
        L \in (1, +\infty] & \rightarrow \exists\lim_{n\rightarrow\infty} a_n = +\infty
    \end{cases}
\end{equation}

\subsubsection{Ordini di infinito}
\begin{equation}
    n^n \ggg n! \ggg q^n (\forall q>1) \ggg n^\alpha (\forall\alpha > 0) \ggg \log_\alpha n (\forall\alpha>1)
\end{equation}


    \section{Funzioni}
\subsection{Domini di funzioni elementari}
\begin{itemize}
    \item $f(x) = \sqrt{x} : x \ge 0$
    \item $f(x) = \log{x} : x > 0$
    \item $f(x) = \frac{1}{x} : x \not = 0$
    \item $f(x) = P(x)^\alpha : x \ge 0$ se $\alpha$ numero illimitato NON periodico.
\end{itemize}

\subsection{Funzioni trigonometriche e iperboliche}
\begin{itemize}
    \item $sin(x)$ e $cos(x)$: $\mathbb{R}$
    \begin{itemize}
        \item $sin^2(x) + cos^2(x) = 1$
        \item $cos(\alpha + \beta) = cos(\alpha)\cdot cos(\beta) + sin(\beta)\cdot sin(\alpha)$
        \item $sin(\alpha + \beta) = sin(\alpha)\cdot cos(\beta) + sin(\beta)\cdot cos(\alpha)$
    \end{itemize}
    \item $tan(x) : \mathbb{R} \backslash \big\{\frac{\pi}{2} + k\pi\big\}$
    \item $arctan(x) : \mathbb{R}$
    \item $sinh(x)$ e $cosh(x)$: $\mathbb{R}$
    \begin{itemize}
        \item $cosh^2(x) - sinh^2(x) = 1$
        \item $cosh(x) = \frac{(e^x + e^{-x})}{2}$
        \item $sinh(x) = \frac{(e^x - e^{-x})}{2}$
        \item $cosh(\alpha + \beta) = cosh(\alpha)\cdot cosh(\beta) + sinh(\alpha)\cdot sinh(\beta)$
        \item $sinh(\alpha + \beta) = sinh(\alpha)\cdot cosh(\beta) + cosh(\alpha)\cdot sinh(\beta)$
    \end{itemize}
\end{itemize}

\subsection{Limiti di funzioni}
\subsubsection{\texorpdfstring{$f$ converge in $l \in \mathbb{R}$}{}}
Sia $f:A\rightarrow\mathbb{R}$ punto di accumulazione per $A$, allora:

\begin{itemize}
    \item $x_0\in\mathbb{R}$:
    \begin{equation}
        \begin{aligned}
            & \exists \lim_{x \to x_0} f(x) = l \quad \text{se:} \\
            & \forall \varepsilon > 0 \qquad \exists \delta > 0 : |f(x) - l| < \varepsilon \qquad \forall x \in A \ \text{con} \ 0 < |x - x_0| < \delta
        \end{aligned}
    \end{equation}
    
    \item $x_0 = \pm\infty$:
    \begin{equation}
        \begin{aligned}
            & \exists \lim_{x \to \pm\infty} f(x) = l \quad \text{se:} \\
            & \forall \varepsilon > 0 \qquad \exists N > 0 : |f(x) - l| < \varepsilon \qquad \forall x \in A \ \text{con} 
            \begin{cases}
                x > N &(+\infty)\\ 
                x < -N\quad&(-\infty)
            \end{cases}
        \end{aligned}
    \end{equation}
\end{itemize}

\subsubsection{\texorpdfstring{$f$ diverge a $\pm\infty$}{}}
Sia $f:A\to\mathbb{R}$ e $x_0$ punto di accumulazione per A, allora:
\begin{itemize}
    \item $x_0\in\mathbb{R}$:
    \begin{equation}
        \begin{aligned}
            & \exists \lim_{x \to x_0} f(x) = \pm\infty \quad \text{se:} \\
            & \forall M > 0 \quad \exists\delta>0:
            \begin{cases}
                f(x)>M &(+\infty)\\ 
                f(x) < -M\ &(-\infty)
            \end{cases} \qquad
            \forall x\in A\ \text{con}\ 0 < |x-x_0| < \delta
        \end{aligned}
    \end{equation}
    
    \item $x_0 = \pm\infty$:
    \begin{equation}
        \begin{aligned}
            & \exists \lim_{x \to \pm\infty} f(x) = \pm\infty \quad \text{se:} \\
            & \forall M > 0 \quad \exists N>0:
            \begin{cases}
                f(x)>M &(+\infty)\\ 
                f(x) < -M\ &(-\infty)
            \end{cases} \qquad
            \forall x\in A\ \text{con}
            \begin{cases}
                x>N &(+\infty)\\ 
                x<-N\ &(-\infty)
            \end{cases}
        \end{aligned}
    \end{equation}
\end{itemize}

\subsection{Funzioni SU, 1-1, biettive}
Sia $f:A\to B$ una funzione.
\begin{itemize}
    \item $f$ si dice \textbf{suriettiva} se:
    \begin{equation}
        f(A) = B
    \end{equation}
    cioè $\forall x\in B \quad \exists x\in A : y = f(x)$. In tal caso scriveremo $f:A\xrightarrow{SU} B$

    \item $f$ si dice \textbf{iniettiva} se:
    \begin{equation}
        f(x) = f(x') \Longrightarrow x = x' \qquad \forall x, x' \in A
    \end{equation}
    cioè $\forall x, x' \in A$ con $x \not = x'$ si ha $f(x) \not = f(x')$. In tal caso scriviamo $f:A\xrightarrow{1-1}B$

    \item $f$ si dice \textbf{biettiva/biunivoca} se $f$ è sia suriettiva sia iniettiva. In tal caso scriviamo $f:A\xrightarrow{SU, 1-1}B$
\end{itemize}

\subsection{Funzioni continue}
Sia $f:A\to\mathbb{R}$ supponiamo $x_0\in A$ e $x_0$ punto di accumulazione per $A$.
La funzione si dice continua in $x_0$ se:
\begin{equation}
    \exists\lim_{x_{x\in A}\to x_0}f(x) = f(x_0)
\end{equation}
\subsubsection{Teorema di Bolzano (o "Teorema degli zeri")}
Sia una funzione $f:A\to\mathbb{R}$, e sia $[a,b]\subseteq A$ un \textbf{intervallo chiuso e limitato}. Supponiamo inoltre che $f \in C([a, b])$, e che assuma agli estremi dell'intervallo valori di segno opposto, cioè: $f(a)\cdot f(b) < 0$.
Allora $f$ ammette almeno uno zero interno ad $[a, b]$, cioè $\exists f(x_0) = 0$.
\subsubsection{Corollario (o "Teorema dei valori interni")}
Sia $I$ intervallo di $\mathbb{R}$, e sia $f:I\to\mathbb{R}$ continua. Il teorema dice che se prendi due punti ($f(x_1)$ e $f(x_2)$) prendo tutti i valori di $y$ che stanno in mezzo tra $f(x_1)$ e $f(x_2)$.
\subsubsection{Teorema di Weierstrass}
Siano $a<b$. Sia $f:[a, b] \to \mathbb{R}$
Supponiamo $f$ continua. Allora Esistono il minimo ed il massimo assoluti per $f$ su $[a, b]$, cioè esistono $x_{min} \text{ e } x_{max} \in [a, b]$.

    \section{Teorema delle restrizioni}
    Sia $f:A\to\mathbb{R}$ e sia $x_0$ punto di accumulazione per $A$. Sia inoltre $A_0 \subseteq A$ e supponiamo che $x_0$ sia punto di accumulazione anche per $A_0$.
    
    Se $\exists \lim_{x_{\in A}\to x_0}f(x)$, allora:
    \begin{equation}
        \exists \lim_{x_{x\in A_0}\to x_0} f(x) = \lim_{x_{x\in A}\to x_0} f(x)
    \end{equation}

    \subsection{Osservazione}
    Se $f:A\rightarrow\mathbb{R}$ con $x_0$ punto di accumulazione per $A$, se trovo:
    \begin{equation}
        \begin{cases}
            A_1 \in A\\ 
            \text{$x_0$ punto di accumulazione per $A_1$}
        \end{cases}
        \qquad
        \text{e}
        \qquad
        \begin{cases}
            A_2 \in A\\ 
            \text{$x_0$ punto di accumulazione per $A_2$}
        \end{cases}
    \end{equation}
    tale che:
    \begin{equation}
        \lim_{x_{x\in A_1}\to x_0} f(x) \not = \lim_{x_{x\in A_2}\to x_0} f(x) \Longrightarrow \not\exists \lim_{x_{x\in A}\to x_0} f(x)
    \end{equation}


    \section{Limiti "assolutamente notevoli"}
    \begin{equation}
        \lim_{x\to 0} \frac{sin(x)}{x} = 1
    \end{equation}
    \begin{equation}
        \lim_{x\to 0} \frac{1- cos(x)}{x^2} = \frac{1}{2}
    \end{equation}
    \begin{equation}
        \lim_{x\to 0} \frac{e^x - 1}{x} = 1
    \end{equation}
    \begin{equation}
        \lim_{x\to 0} \frac{log(x+1)}{x} = 1
    \end{equation}

    \section{Teorema del cambaimento di variabile}
    Siano $A, B \subseteq \mathbb{R}$ e sia $h:A\to B$ e $g:B\to\mathbb{R}$. Supponiamo $x_0$ punto di accumulazione per $A$ e che:
    \begin{equation}
        \exists\lim_{x\to x_0} h(x) = l \qquad (l\in\mathbb{R} \text{ o } \pm\infty)
    \end{equation}
    Supponiamo che che $l$ sia di accumulazione per $B$ e che:
    \begin{equation}
        \exists\lim_{y\to l}g(y) = m
    \end{equation}
    Allora, se vale uno dei seguenti due casi:
    \begin{itemize}
        \item $l\not\in B$
        \item $l\in B \text{ e } g(l) = m$
    \end{itemize}
    risulta che:
    \begin{equation}
        \exists\lim_{n\to x_0} g(h(x)) = m
    \end{equation}


    % ***************************************
    % SECONDO PARZIALE
    % ***************************************

    \section{Formule trigonometriche}
    \begin{equation}
        sin(2a) = 2sin(x)cos(x)
    \end{equation}
    \begin{equation}
        tan(2a) = \frac{2tan(a)}{1-tan^2(a)}
    \end{equation}

    \section{Derivate}
Sia $f: A\to\mathbb{R}$, e sia $x_0$ punto interno ad A. Diciamo che $f$ è derivabile in $x_0$ se:
\begin{equation}
    \exists\lim_{x\to x_0}\frac{f(x)-f(x_0)}{x-x_0} = l \in\mathbb{R}
\end{equation}

% \subsection{Derivata di una retta}
% Sia $l:\mathbb{R}\to\mathbb{R}$ con $l(x) = mx+q$, pertanto:
% \begin{equation}
%     g(x)=\frac{l(x)-l(x_0)}{x-x_0} = \frac{(mx+q) - (mx_0+q)}{x-x_0} = m
% \end{equation}
% Allora:
% \begin{equation}
%     \exists\lim_{x\to x_0} g(x) = m
% \end{equation}
% In particolare \textbf{le funzioni costanti hanno derivata nulla}.

\subsection{Retta tangente}
\begin{equation}
    y = f'(x_0)(x-x_0) + f(x_0)
\end{equation}

\subsection{Regole di derivazione}
\begin{equation}
    (f+g)'(x_0) = f'(x_0) + g'(x_0)
\end{equation}
\begin{equation}
    (f\cdot g)'(x_0) = f'(x_0) g(x_0) + f(x_0)g'(x_0)
\end{equation}
\begin{equation}
    \bigg(\frac{f}{g}\bigg)'(x_0) = \frac{f'(x_0) g(x_0) - f(x_0)g'(x_0)}{g^2(x_0)} \qquad \bigg(\frac{1}{g}\bigg)'(x_0) = \frac{-g'(x_0)}{g^2(x_0)}\end{equation}
\begin{equation}
    (g\ o\ f)'(x_0) = g'(f(x_0))\cdot f'(x_0)
\end{equation}
\begin{equation}
    (f^{-1})'(y_0) = \frac{1}{f'(x_0)}
\end{equation}

\subsection{Derivate notevoli}
Calcolo le derivate della funzione $f$ in $x_0$ con:
\begin{equation}
    f'(x_0) = \lim_{t \to 0}\frac{f(x_0 + t)- f(x_0)}{t}
\end{equation}
Allora:
\begin{equation}
    f(x) = x^n \qquad f'(x) = nx^{n-1}
\end{equation}
\begin{equation}
    f(x) = x^{-n} = \frac{1}{x^n} \qquad f'(x) = -nx^{-n-1} = \frac{-n}{x^{n+1}}
\end{equation}
\begin{equation}
    f(x) = e^x \qquad f'(x) = e^x
\end{equation}
\begin{equation}
    f(x) = e^{-x} \qquad f'(x) = -e^{-x}
\end{equation}
\begin{equation}
    f(x) = a^x \qquad f'(x) = log(a) \cdot a^x
\end{equation}
\begin{equation}
    f(x) = log(x) \qquad f'(x) = \frac{1}{x}
\end{equation}
\begin{equation}
    f(x) = x^\alpha \qquad f'(x) = \alpha x^{\alpha - 1} 
\end{equation}
\begin{equation}
    f(x) = \sqrt{x} = x^\frac{1}{2} \qquad f'(x) = \frac{1}{2\sqrt{x}}
\end{equation}
\begin{equation}
    f(x) = sin(x) \qquad f'(x) = cos(x)
\end{equation}
\begin{equation}
    f(x) = cos(x) \qquad f'(x) = -sin(x)
\end{equation}
\begin{equation}
    f(x) = tan(x) = \frac{sin(x)}{cos(x)} \qquad f'(x) = 1 + tan^2(x)
\end{equation}
\begin{equation}
    f(x) = arcsin(x) \qquad f'(x) = \frac{1}{\sqrt{1-x^2}}
\end{equation}
\begin{equation}
    f(x) = arcos(x) \qquad f'(x) = \frac{-1}{\sqrt{1-x^2}}
\end{equation}
\begin{equation}
    f(x) = arctan(x) \qquad f'(x) = \frac{1}{1+x^2}
\end{equation}

Nel caso la $x$ si presentasse sia alla base che all'esponente (es. $f(x) = x^x$, $f(x) = (sin(x))^{cos(x)}$) potremmo riscrivere la funzione come:  $f(x) = e^{log(x^x)} = e^{xlog(x)}$.
Sviluppando la derivata avremmo allora: $f'(x) = e^{xlog(x)}\cdot \Big[log(x) + x\cdot\frac{1}{x}\Big] = e^{xlog(x)}\cdot [log(x) + 1]$. Ma essendo $x^x = e^{log(x^x)}$ possiamo riscrivere il risultato come: $f'(x) = x^x \cdot [log(x)+1]$.
    
    \section{Studio di funzione}

\begin{enumerate}
    \item \textbf{Dominio}
    \item \textbf{Simmetrie}
          \begin{itemize}
            \item Pari: $f(-x) = f(x)$
            \item Dispari: $f(-x) = -f(x)$
          \end{itemize}
    \item \textbf{Studio del segno e intersezioni}
          \begin{itemize}
            \item Maggiore: $f(x) > 0$
            \item Intersezioni con $x$: $f(x) = 0$
            \item Intersezioni con $y$: $f(0)$
          \end{itemize}
    \item \textbf{Limiti e asintoti}\\
          I limiti vanno calcolati nei \textit{punti di accumulazione} e agli \textit{estremi del dominio}.
    \item \textbf{Studio del segno della derivata}
          \begin{itemize}
            \item Derivata positiva -> funzione crescente
            \item Derivata negativa -> funzione decrescente
          \end{itemize}
\end{enumerate}


    
    \section{Integrali}
\begin{equation}
    \int_a^bf(x)dx
\end{equation}
Restituisce l'area (con segno) delle regione di piano compresa tra il grafico di $f(x)$, l'asse $x$ e le rette verticali $x=a$ e $x=b$.
Se invece $g(x) = K$ $\forall x\in[a, b]$, allora l'integrale di $g$ sull'intervallo è $\int_a^b g(x)dx = (b-a)\cdot K$.

\subsection{Come calcolarli}
Per calcolare $\int_a^bf(x)dx$ devo:
\begin{enumerate}
    \item trovare una funzione che, nell'intervallo $[a, b]$, abbia $f(x)$ come derivata (primitiva di $f$).
    \item la calcolo negli estremi della zona di integrazione (calcolo $F(b)$ e $F(a)$).
    \item sottraggo i due valori: $\int_a^bf(x)dx = F(b)-F(a)$
\end{enumerate}

\subsection{Primitive}
Sia $I$ intervallo di $\mathbb{R}$, e sia $f: I\to\mathbb{R}$. Diciamo che: $$F: \mathring{I}\to\mathbb{R}$$
è primitiva di $f$ se:
\begin{itemize}
    \item $F$ è derivabile $\forall x \in \mathring{I}$
    \item $F'(x)=f(x) \qquad\forall x \in \mathring{I}$
\end{itemize}
\[
\renewcommand{\arraystretch}{1.5} % Aumenta lo spazio verticale tra le righe
\begin{array}{|c|c|}
\hline
\mathbf{f(x)} & \mathbf{F(x)} \\ \hline
\cos(x) & \sin(x) \\ \hline
\sin(x) & -\cos(x) \\ \hline
\tan(x) & -\ln\lvert\cos(x)\lvert \\ \hline
\arccos(x) & x \arccos(x) - \sqrt{1 - x^2} \\ \hline
\arcsin(x) & x \arcsin(x) + \sqrt{1 - x^2} \\ \hline
e^{ax} & \frac{e^{ax}}{a} \\ \hline
\alpha^x & \frac{\alpha^x}{\ln(\alpha)} \\ \hline
x^n & \frac{x^{n+1}}{n+1} \quad \forall n \neq -1 \\ \hline
\frac{1}{x} & \ln(\lvert x \rvert) \\ \hline
\frac{1}{1+x^2} & \arctan(x) \\ \hline
\end{array}
\]


\subsection{Proprietà}
\begin{equation}
    \int[f(x)\pm g(x)]dx = \int f(x)dx \pm \int g(x)dx
\end{equation}
\begin{equation}
    \int Kf(x)dx = K\int f(x)dx
\end{equation}
\begin{equation}
    \int_a^bf(x)dx = \int_a^cf(x)dx  + \int_c^bf(x)dx
\end{equation}


\subsection{Metodi}
\subsubsection{Integrazione per parti}
\begin{equation}
    \int f(x)g'(x) = f(x)g(x) - \int f'(x)g(x)dx
\end{equation}

\end{document}
